\documentclass[11pt, a4paper]{article}

\usepackage{amsmath, amssymb}
\usepackage[margin=1in]{geometry}
\usepackage{graphicx}
\usepackage{hyperref}
\usepackage{enumitem}
\usepackage{xcolor}
\usepackage{booktabs}
\usepackage{tikz}
\usetikzlibrary{arrows.meta, positioning}

\hypersetup{colorlinks=true, linkcolor=blue!50!black, citecolor=blue!50!black, urlcolor=blue!50!black}

% --- draft annotations ---------------------------------------------
\newcommand{\note}[1]{\textcolor{red!70!black}{\textit{[#1]}}}
\newcommand{\todo}[1]{\textcolor{blue!70!black}{\textit{[TODO: #1]}}}

% --- shorthand ------------------------------------------------------
\newcommand{\R}{\mathbb{R}}
\newcommand{\Rpos}{\mathbb{R}_{\geq 0}}
\newcommand{\rstar}{r^{\star}}
\newcommand{\Sstar}{S^{\star}}
\newcommand{\Estar}{\mathcal{E}^{\star}}
\newcommand{\Gstar}{G^{\star}}
\newcommand{\Astar}{A_{\star}}
\newcommand{\Bstar}{B_{\star}}
\newcommand{\Dstar}{D_{\star}}

\title{Identifiability of co-transcriptional splicing rates from steady-state isoform distributions}
\author{Swapnil Mittal \and Justin Kinney}
\date{Working draft --- \today}

\begin{document}
\maketitle

\begin{abstract}
We study when kinetic rates of a co-transcriptional splicing system can be recovered from a single steady-state observation of its isoform proportion distribution (IPD). The system is modelled as a continuous-time Markov chain on a directed acyclic graph (the DADA graph) whose nodes are partial-splicing intermediates and whose edges fire with unknown rates. We parametrise families of rate vectors by a low-dimensional splicing regime vector and ask, for each region of regime space, whether the inverse map from IPD to rates is well-posed. Recovery from a single IPD turns out to be possible only in trivial cases (one intron) or under restrictive structural assumptions (Kinetic regime, where all rates within a junction are equal). For richer regimes, two distinct levels of structural information restore identifiability: knowing the active subgraph, or observing the same gene under multiple export contexts.
\end{abstract}

\section{Introduction}

\note{Motivate co-transcriptional kinetic competition and why \emph{rates} matter beyond outcome frequencies. Contrast with sequence-based alternative-splicing predictors. State the gap that motivates the inverse-problem framing.}

The aim of this note is conceptual, not algorithmic. We do not propose a new estimator. We characterise the well-posedness of an inverse problem that several existing methods (NNLS-based and otherwise) tacitly assume to be solvable, and we work out what additional information --- structural priors, multiple observations of the same gene --- is needed to make the problem well-posed in cases where it is not.

\section{The DADA graph}

\subsection{Splice sites, junctions, and compatibility}

A pre-mRNA carries an alternating sequence of donor (\textsf{D}) and acceptor (\textsf{A}) splice sites. We write the gene's site sequence as a string in the alphabet $\{\textsf{D}, \textsf{A}\}$ --- the \emph{DADA string}. A \emph{junction} is an ordered pair $j = (d, a)$ of a donor at site $d$ and a later acceptor at site $a > d$; junction $j$ corresponds to the molecular event of splicing out the intron between sites $d$ and $a$.

Two junctions $j = (d, a)$ and $j' = (d', a')$ are \emph{compatible} if (a) they share no splice site --- $d \neq d'$ and $a \neq a'$ --- and (b) their open intervals $(d, a)$ and $(d', a')$ are either disjoint or properly nested.

\subsection{States, edges, and isoforms}

Define the state set $\mathcal{V}$ as the collection of \emph{antichains} of compatible junctions: each $S \in \mathcal{V}$ is a set of mutually compatible junctions. The biological reading is that $S$ records exactly the splices that have already been made; the molecule in state $S$ is the partial isoform with those introns removed. The state $\emptyset$ is the unspliced pre-mRNA; the maximal antichains are mature isoforms.

The DADA graph is
\[
G \;=\; (\mathcal{V}, \mathcal{E}), \qquad \mathcal{E} \;=\; \big\{\, (S, S \cup \{j\}) \,:\, S \in \mathcal{V},\ j \notin S,\ j \text{ compatible with all of } S \,\big\}.
\]
Each edge fires a single junction. Edges are written $e = S \to S \cup \{j\}$, and we say edge $e$ is \emph{labelled by} junction $j(e) = j$; multiple edges may share a junction label, since the same splicing reaction can fire from different intermediate isoforms. The rate of edge $e$ is $r_e \geq 0$. The graph is rooted at $\emptyset$ and is, by construction, a directed acyclic graph (DAG): each edge strictly enlarges the set of completed splices.

\subsection{Worked example: \texttt{DADA} (\texorpdfstring{$K = 2$ introns}{K=2 introns})}
\label{sec:dada-example}

For the DADA string \texttt{DADA} (sites $\textsf{D}_0\,\textsf{A}_1\,\textsf{D}_2\,\textsf{A}_3$) the junctions are
\[
j_1 = (0, 1), \qquad j_2 = (0, 3), \qquad j_3 = (2, 3).
\]
Compatibility:
\begin{itemize}[noitemsep, topsep=2pt]
  \item $j_1$ and $j_2$ share donor $\textsf{D}_0$: incompatible.
  \item $j_2$ and $j_3$ share acceptor $\textsf{A}_3$: incompatible.
  \item $j_1 = (0,1)$ and $j_3 = (2,3)$: disjoint intervals: compatible.
\end{itemize}
The state set $\mathcal{V} = \{\,\emptyset,\ \{j_1\},\ \{j_2\},\ \{j_3\},\ \{j_1, j_3\}\,\}$ has $N = 5$ elements; the graph has $M = 5$ edges:
\begin{align*}
e_1: \emptyset \to \{j_1\},\ j(e_1) = j_1, \quad &
e_2: \emptyset \to \{j_2\},\ j(e_2) = j_2, \\
e_3: \emptyset \to \{j_3\},\ j(e_3) = j_3, \quad &
e_4: \{j_1\} \to \{j_1, j_3\},\ j(e_4) = j_3, \\
e_5: \{j_3\} \to \{j_1, j_3\},\ j(e_5) = j_1. &
\end{align*}
Note that junctions $j_1$ and $j_3$ each label two edges (different intermediate states from which the same splicing reaction can fire), while $j_2$ labels only one edge because it is incompatible with everything else.

\todo{Insert figure: DADA state graph with nodes $\emptyset, \{j_1\}, \{j_2\}, \{j_3\}, \{j_1,j_3\}$, edges labelled by $j(e)$.}

\section{The forward model}
\label{sec:forward}

Each non-root state $i$ is exported (transcribed off the chromatin and observed as mature mRNA) at a known rate $\lambda_i \geq 0$. Pre-mRNA enters the system at the root $i = 0$ at unit production rate. The probability $p_i$ of finding a single molecule in state $i$ at steady state satisfies, for each $i \in \mathcal{V}$,
\begin{equation}
\underbrace{\sum_{e:\ \mathrm{dst}(e) = i} r_e\, p_{\mathrm{src}(e)}}_{\text{flux in}}
\;-\;
\underbrace{\sum_{e:\ \mathrm{src}(e) = i} r_e\, p_i}_{\text{flux out}}
\;-\;
\lambda_i\, p_i
\;=\;
-\delta_{i, 0}.
\label{eq:balance-row}
\end{equation}
The root row is dependent on the others (mass conservation: total inflow = total export), so we drop it and stack the remaining $N - 1$ equations into matrix form,
\begin{equation}
A(p)\, r \;=\; c(p), \qquad A \in \R^{(N - 1) \times M}, \qquad c \in \R^{N - 1},
\label{eq:balance}
\end{equation}
with
\[
A_{i, e} \;=\; \begin{cases} +\,p_{\mathrm{src}(e)} & \text{if } \mathrm{dst}(e) = i, \\ -\,p_{\mathrm{src}(e)} & \text{if } \mathrm{src}(e) = i, \\ 0 & \text{otherwise}, \end{cases}
\qquad
c_i \;=\; \lambda_i\, p_i.
\]
Equation~\eqref{eq:balance} is the central object. We observe $p$ (the IPD), $\lambda$ is known, so $A$ and $c$ are computable; the rates $r$ are the unknowns.

\section{The inverse problem}
\label{sec:inverse}

Given the IPD $p$ and the export rates $\lambda$, we want to recover the rate vector $r$. A natural formulation is
\begin{equation}
\widehat{r} \;=\; \arg\min_{r \geq 0} \; \big\| A(p)\, r - c(p) \big\|_2^2.
\label{eq:nnls}
\end{equation}
The non-negativity constraint is biologically essential: rates are physical reaction frequencies. We refer to the optimisation problem~\eqref{eq:nnls} as the \emph{inverse problem}; the question of this paper is whether $\widehat{r}$ equals the truth $\rstar$, and under what conditions.

\section{The splicing regime vector}
\label{sec:srv}

\subsection{Generative model for rates}

We model the truth $\rstar$ as a draw from a parametric family controlled by a low-dimensional regime vector
\[
\theta \;=\; (\eta,\ \sigma,\ \alpha) \;\in\; [0, 1] \times \R_{\geq 0} \times \R_{> 0}.
\]
The Mixed-Dirichlet generator at $\theta$ samples a rate vector by:
\begin{enumerate}[leftmargin=*, itemsep=2pt, topsep=2pt]
  \item For each junction $j$, draw an amplitude $\rho_j \sim \mathrm{LogNormal}(\log R_0,\ \sigma^2)$.
  \item For each non-root state $i$ with $m_i$ incoming edges, set the active cardinality
    \[
    k_i \;=\; \max\bigl(1,\ \mathrm{round}(\eta\,m_i)\bigr),
    \]
    pick a uniform random size-$k_i$ subset $S_i$ of the $m_i$ incoming edges, and draw weights $w \sim \mathrm{Dirichlet}(\alpha, \dots, \alpha)$ over $S_i$.
  \item Set $r_e = \rho_{j(e)} \cdot k_i \cdot w_e$ if $e \in S_i$ (where $i = \mathrm{dst}(e)$), else $r_e = 0$.
\end{enumerate}
The active mask is $\Sstar = \{e : r_e > 0\}$ and the active subgraph $\Gstar$ is the subgraph induced on edges in $\Sstar$.

\subsection{Four canonical regimes}

The corners and faces of regime space have biologically natural interpretations:
\begin{itemize}[leftmargin=*, itemsep=2pt, topsep=2pt]
  \item \textbf{Null}: $\sigma = 0$, $\alpha \to \infty$. Degenerate: all rates equal $R_0$. Useful as a baseline.
  \item \textbf{Predetermined}: $\eta = 0$. Each child has exactly one active parent, so $\Gstar$ is a tree (a spanning tree of the active state subgraph).
  \item \textbf{Kinetic}: $\eta = 1$, $\alpha \to \infty$. Every edge is active and the Dirichlet weights are uniform within each junction, so $r_e = \rho_{j(e)}$ for every active edge of junction $j$.
  \item \textbf{Mixed}: $\eta \in (0, 1)$, $\alpha$ moderate. Partial branching with uneven weights. The biologically realistic regime, and the hard one for inverse inference.
\end{itemize}
The four are not exhaustive; they are landmarks. The whole prior cube $\theta \in [0,1] \times [0, \infty) \times (0, \infty)$ is a continuous interpolation between them.

\section{What is identifiable, and when}
\label{sec:identifiability}

We work throughout this section with a single noiseless IPD $p$ and known export rates $\lambda$.

\subsection{The relevant linear-algebraic objects}

Given an active edge set $\Estar$ define the active subgraph $\Gstar = (\mathcal{V}^{\star}, \Estar)$ where $\mathcal{V}^{\star}$ is the set of states touched by edges in $\Estar$ together with the root; let $N^{\star} = |\mathcal{V}^{\star}|$ and $M^{\star} = |\Estar|$. We assume throughout that $\Gstar$ is connected (always true for our DADA graphs because of the rooted construction).

The masked balance matrix $\Astar$ is the column restriction of $A$ in~\eqref{eq:balance} to columns indexed by $\Estar$. It admits the factorisation
\[
\Astar \;=\; \Bstar \cdot \Dstar,
\]
where $\Bstar$ is the (root-removed) signed incidence matrix of $\Gstar$ --- with $(\Bstar)_{i, e} = +1$ if $\mathrm{dst}(e) = i$, $-1$ if $\mathrm{src}(e) = i$, and $0$ otherwise --- and $\Dstar$ is the diagonal matrix with $(\Dstar)_{e,e} = p_{\mathrm{src}(e)}$. When all $p_{\mathrm{src}(e)} > 0$, $\Dstar$ is invertible and
\[
\mathrm{rank}(\Astar) \;=\; \mathrm{rank}(\Bstar) \;=\; N^{\star} - 1.
\]
This last equality is the standard fact that the (reduced) signed incidence matrix of a connected graph on $N^{\star}$ vertices has rank $N^{\star} - 1$.

The kernel of $\Astar$ is therefore
\[
\ker(\Astar) \;=\; \Dstar^{-1} \ker(\Bstar) \;\cong\; \ker(\Bstar),
\]
which is the cycle space of $\Gstar$ (treated as an undirected graph for cycle-space purposes), of dimension equal to the cyclomatic number
\begin{equation}
c(\Gstar) \;=\; M^{\star} - N^{\star} + 1.
\label{eq:cyclomatic}
\end{equation}

A non-trivial element of $\ker(\Bstar)$ corresponds to a cycle in the underlying undirected graph of $\Gstar$ --- equivalently, two distinct directed paths between the same pair of states. Although $\Gstar$ has no directed cycles (it is a DAG), the existence of multiple parallel directed paths between the same pair of vertices means rate flux can be \emph{redistributed} between alternatives without changing the IPD.

\subsection{Single IPD with known active subgraph}
\label{sec:single-ipd-mask}

If we are given the exact active set $\Estar$, the inverse problem reduces to solving \eqref{eq:nnls} restricted to columns of $\Estar$. The picture splits cleanly on $c(\Gstar)$:

\paragraph{Tree case ($c(\Gstar) = 0$).} Then $M^{\star} = N^{\star} - 1 = \mathrm{rank}(\Astar)$, so $\Astar$ has full column rank, the system $\Astar r = c$ has at most one solution, the truth $\rstar$ is one such, and NNLS returns $\rstar$ exactly.

\paragraph{Cyclic case ($c(\Gstar) > 0$).} Then $\ker(\Astar)$ has positive dimension. For a generic truth $\rstar$ in the strictly positive orthant, $\rstar$ is in the relative interior of the polytope
\[
\big\{\, r \geq 0 \,:\, \Astar r = c \,\big\},
\]
which has dimension $c(\Gstar) > 0$. Multiple non-negative rate vectors then produce the same IPD, and NNLS alone (without further structure) cannot distinguish them. The corner of the polytope picked by Lawson--Hanson NNLS depends on algorithmic details rather than on the data.

\paragraph{Numerical aside.} Even in the tree case, plain NNLS can be numerically poor on graphs of large depth because $\Astar = \Bstar \Dstar$ inherits a large condition number when source-state probabilities span many orders of magnitude. We have observed that on the tree subgraph of \texttt{DADADADADADA} ($K = 6$ introns), per-edge NNLS frequently fails despite full column rank. The augmented system used by junction-shrinkage NNLS (Section~\ref{sec:junction-shrinkage}) regularises the conditioning sufficiently to recover the truth at machine precision.

\subsection{Single IPD without the mask}
\label{sec:no-mask}

In the realistic case where we observe only $p$ and the topology $G$ (not which edges are active), the situation is strictly worse. The full balance matrix $A$ has $M$ columns and $N - 1$ rows with $M > N - 1$ for every non-trivial DADA graph; therefore
\[
\dim \ker(A) \;\geq\; M - (N - 1) \;>\; 0,
\]
and the NNLS solution set has positive dimension regardless of $c(\Gstar)$. The truth $\rstar$ is one of many non-negative rate vectors fitting the same $p$.

A separate issue is that distinct active subgraphs can produce identical IPDs. The \emph{cardinality shell} of a realisation is the set of active masks consistent with the cardinality profile $(k_i)_{i \in \mathcal{V}^{\star}}$. We have observed empirically (\note{on \texttt{DADADA}, $K = 3$, with $\eta = 0.5$, all 96 elements of the shell admit non-negative rate vectors fitting the same IPD to floating-point precision}). Cardinality is identifiable from the IPD; mask identity, generically, is not.

\subsection{Junction-shrinkage and the Kinetic regime}
\label{sec:junction-shrinkage}

Suppose we additionally constrain the rates to lie on the \emph{junction-equal subspace}
\[
W_{\mathrm{eq}} \;=\; \big\{\, r \in \R^{M^{\star}} \,:\, r_{e} = r_{e'} \text{ whenever } j(e) = j(e') \,\big\}.
\]
Inside $W_{\mathrm{eq}}$, the unknowns collapse from $M^{\star}$ edge rates to $J^{\star}$ junction rates, where $J^{\star}$ is the number of junctions touched by edges in $\Estar$. The pooled balance matrix $A_{\mathrm{pool}}$ then has shape $(N - 1) \times J^{\star}$, and for typical DADA topologies $J^{\star} \ll N - 1$, so $A_{\mathrm{pool}}$ is generically full column rank.

Two consequences:

\begin{enumerate}[leftmargin=*, itemsep=2pt, topsep=2pt]
  \item Pooled NNLS (which solves~\eqref{eq:nnls} on $A_{\mathrm{pool}}$ rather than $A$) recovers the junction rates uniquely from a single IPD. \emph{No active mask is required}: an inactive edge contributes zero to the column-aggregation in its junction. This is the Kinetic regime ($\eta = 1$, $\alpha \to \infty$), because the truth $\rstar$ then satisfies $\rstar \in W_{\mathrm{eq}}$ exactly --- all active edges of a junction share the rate $\rho_j$.

  \item The junction-shrinkage NNLS solver augments the balance equations with Helmert contrasts, equivalent to a soft penalty for deviation from $W_{\mathrm{eq}}$. The augmented matrix is overdetermined regardless of $c(\Gstar)$. When $\rstar \in W_{\mathrm{eq}}$, the augmented system is consistent and recovers $\rstar$ exactly (this includes both the Predetermined regime, where every junction has at most one active edge so the constraints are vacuous, and the Kinetic regime). When $\rstar \notin W_{\mathrm{eq}}$, the augmented system is inconsistent and the LS solution is biased toward $W_{\mathrm{eq}}$, giving $\widehat{r} \neq \rstar$.
\end{enumerate}

The boundary between "junction-shrinkage works" and "junction-shrinkage fails" is therefore not the cyclomatic boundary $c(\Gstar) = 0$, but the inclusion $\rstar \in W_{\mathrm{eq}}$. Under the Mixed-Dirichlet generator, $\Pr(\rstar \in W_{\mathrm{eq}}) \to 1$ as $\alpha \to \infty$ (the Kinetic limit) and $\Pr(\rstar \in W_{\mathrm{eq}}) = 1$ when $\eta = 0$ (Predetermined, vacuous constraints). For Mixed regimes at finite $\alpha$, the probability is generically zero.

\subsection{Worked example: \texorpdfstring{$K = 2$ (\texttt{DADA})}{K=2 (DADA)}}

For \texttt{DADA} (Section~\ref{sec:dada-example}) the full graph has $M = 5$, $N = 5$, and
\[
c(G) \;=\; M - N + 1 \;=\; 1.
\]
The single cycle is the diamond at state $\{j_1, j_3\}$, reachable through $\{j_1\} \xrightarrow{e_4} \{j_1, j_3\}$ or through $\{j_3\} \xrightarrow{e_5} \{j_1, j_3\}$. Concrete consequences:

\begin{itemize}[leftmargin=*, itemsep=2pt, topsep=2pt]
  \item \textbf{Predetermined, or Mixed at $\eta < 0.75$.} The branching state $\{j_1, j_3\}$ has $m = 2$ parents, and the active cardinality $k = \max(1, \mathrm{round}(2\eta))$ equals $1$ for $\eta < 0.75$. Only one of $\{e_4, e_5\}$ is active. The active subgraph is a $4$-edge tree ($c(\Gstar) = 0$). Identifiable \emph{given the mask}, but not from the IPD alone --- the cardinality shell of size $2$ (active edge $e_4$ \emph{or} $e_5$) is the source of ambiguity.

  \item \textbf{Kinetic ($\eta = 1$, $\alpha \to \infty$).} All $5$ edges active, $\rstar \in W_{\mathrm{eq}}$, pooled NNLS solves a $4 \times 3$ system with unique solution. Identifiable from the IPD alone, no mask required.
\end{itemize}

The trivial graph \texttt{DA} ($K = 1$, $N = 2$, $M = 1$, $c = 0$) is identifiable under any prior because the single rate is read off from the single non-root balance equation directly.

\section{Paths to enlarge identifiability}
\label{sec:paths}

Three structural augmentations restore identifiability outside the cases of Section~\ref{sec:identifiability}. We describe each briefly; each suggests a different experimental or computational direction.

\subsection{Multiple export contexts}

\note{Notation: we use $K$ for the number of introns elsewhere; here we use $T$ (number of tissues / contexts) to avoid the clash.}

Suppose we observe the same gene's IPD under $T$ distinct export profiles $\lambda^{(1)}, \dots, \lambda^{(T)}$ --- e.g.\ the same gene measured in different cell types or conditions, in which the underlying junction rates $r$ are conserved but per-state decay rates differ. Stacking the $T$ balance systems yields
\[
\underbrace{\begin{bmatrix} A^{(1)}(p^{(1)}) \\ \vdots \\ A^{(T)}(p^{(T)}) \end{bmatrix}}_{\text{$T(N-1) \times M^{\star}$}} r \;=\; \begin{bmatrix} c^{(1)} \\ \vdots \\ c^{(T)} \end{bmatrix}.
\]
Each block alone has rank $N^{\star} - 1$, but the column scalings $\Dstar^{(t)}$ differ across blocks because the $p^{(t)}$ differ; the kernels $\ker(A^{(t)}_{\star})$ are accordingly different "rotations" of the same underlying cycle space $\ker(\Bstar)$. Generically, $T \geq \lceil M^{\star} / (N^{\star} - 1) \rceil$ contexts suffice to make the stacked matrix full column rank, recovering $\rstar$ exactly without knowing the active mask. For most practical DADA topologies this number is small (2 or 3).

\note{We should validate this dimension bound on real data --- to what extent is GTEx- or ENCODE-level cell-type variation \emph{generic} in the sense required?}

\subsection{Sequence-based priors on active junctions}

Pre-trained sequence models (SpliceAI, AlphaGenome) predict per-junction probabilities of splice-site usage from sequence alone. Thresholding these declares some junctions inactive, eliminating their edges from $\mathcal{E}$ and shrinking $M^{\star}$. A prior strong enough to bring $M^{\star}$ down to $N^{\star} - 1$ recovers the tree case directly. Even partial reductions help: each removed edge can decrease $c(\Gstar)$ by $0$ or $1$, and the cardinality shell shrinks correspondingly.

\subsection{Direct path observation (long-read)}

Long-read sequencing observes individual transcript paths through the DADA graph rather than aggregate isoform proportions. Empirical edge frequencies from many reads give a direct estimate of the active mask $\Sstar$. Combined with NNLS on the masked balance system, this is the cleanest route in current data modalities and arguably the most under-used.

\section{Summary}

The state of identifiability is summarised in Table~\ref{tab:summary}.

\begin{table}[h]
\centering
\renewcommand{\arraystretch}{1.3}
\begin{tabular}{lcccc}
\toprule
& \multicolumn{2}{c}{\textbf{Single IPD}} & \multicolumn{2}{c}{\textbf{Multi-context IPDs}} \\
\cline{2-3} \cline{4-5}
\textbf{Splice regime} & without mask & with mask & without mask & with mask \\
\midrule
Trivial ($K = 1$, \texttt{DA}) & yes & yes & yes & yes \\
Predetermined / Mixed ($K \geq 2$) & no & yes (tree) & yes \note{check} & yes \\
Kinetic & yes (junction-equal) & yes & yes & yes \\
\bottomrule
\end{tabular}
\caption{Identifiability of the rate vector $\rstar$ as a function of the splice regime and the structural information available. "Tree" means the active subgraph $\Gstar$ has cyclomatic number zero; "junction-equal" means the truth lies in $W_{\mathrm{eq}}$. The mark "\note{check}" indicates an empirical claim we have not yet rigorously verified.}
\label{tab:summary}
\end{table}

The single most important takeaway: \emph{from a single IPD with no structural prior, the only splice regime that admits identifiable rates is Kinetic}, because that is the only regime where the truth's structure (junction-equal) collapses the $M^{\star}$-dimensional inverse problem to the lower-dimensional pooled problem on $J^{\star}$ junction rates. Every other regime requires either the active mask, or multiple observations from different export contexts, or a sequence-based prior that shrinks the candidate active edge set.

\section{Open questions}

\begin{itemize}[leftmargin=*, itemsep=4pt, topsep=2pt]
  \item How non-generic are real cell-type-paired export profiles? Does GTEx- or ENCODE-level across-tissue variation suffice to satisfy the multi-context recovery condition for typical genes?
  \item How sharp is the "junction-equal" boundary for biological data? In real splicing, within-junction weight ratios are presumably non-uniform; how close to equal is "equal enough" for junction-shrinkage NNLS to recover useful (if biased) rates?
  \item What is the right way to combine sequence-based junction priors with the multi-context approach to handle the largest practical DADA graphs ($K \geq 8$ introns), where exhaustive enumeration of cardinality shells or active subgraphs is infeasible?
  \item Are there genes whose splicing structure reliably falls in the Kinetic regime in some tissues? For those, single-IPD inversion via pooled NNLS is exact, with no auxiliary information needed --- a directly testable prediction.
\end{itemize}

% --------------------------------------------------------------
\bibliographystyle{plain}
\begin{thebibliography}{9}

\bibitem{slawski-hein-2013}
Martin Slawski and Matthias Hein.
\newblock Non-negative least squares for high-dimensional linear models: Consistency and sparse recovery without regularisation.
\newblock \emph{Electronic Journal of Statistics}, 7:3004--3056, 2013.

\bibitem{churchman-2020}
\note{nano-COP / Churchman lab paper --- fill in correct reference.}

\bibitem{drexler-2020}
\note{Drexler et al.\ 2020 on K562 chromatin nano-COP --- fill in.}

\bibitem{alphagenome}
\note{Avsec et al., AlphaGenome --- fill in.}

\bibitem{biggs-1993}
Norman Biggs.
\newblock \emph{Algebraic Graph Theory}.
\newblock Cambridge University Press, 2nd edition, 1993.
\newblock (Standard reference for cycle space, signed incidence matrix rank, and the cyclomatic number.)

\end{thebibliography}

\end{document}
